\section{Polyhedral degeneration and quantum volume asymptotics}
\label{sec:quantumdegeneration}

\subsection{Three different geometric quantities}

The connection to the volume conjecture becomes precise only after
separating three quantities:
\[
 A=\operatorname{Area}(\partial P),\qquad
 W(P)=\operatorname{Vol}_{\HH^3}(P),\qquad
 V_{\WP}=\int_{\M}\frac{\omega_{\WP}^d}{d!}.
\]
The first is a surface area, the second an enclosed three-dimensional
volume, and the third a measure of a space of shapes. The ordinary
knot volume conjecture concerns three-dimensional volume. For a
hyperbolic knot, in the reduced colored-Jones normalization
$J_N(\text{unknot};q)=1$, its familiar form is
\begin{equation}
 \lim_{N\to\infty}\frac{2\pi}{N}
 \log\left|J_N(K;e^{2\pi i/N})\right|
 =\operatorname{Vol}(S^3\setminus K).
 \label{eq:volumeconjecture}
\end{equation}
This remains a conjecture in general; see the colored-Jones
formulation in~\cite{MurakamiMurakami} and the tetrahedral
results in~\cite{Costantino}.

Suppose a family of compact convex polyhedra $P_A$ has hyperbolic
boundary area $A\downarrow0$. Make the geometric convergence
precise by choosing isometries into a Klein chart for which the
rescaled convex bodies $A^{-1/2}P_A$ converge in Euclidean
Hausdorff distance to a \emph{nondegenerate} polyhedron
$P_{\mathrm E}$ of surface area one. The rescaled bodies then lie
in a common bounded set, the hyperbolic volume density tends
uniformly to the Euclidean density, and Euclidean volume is
continuous under this convex-body convergence. Hence
\begin{equation}
 W(P_A)=A^{3/2}W_{\mathrm E}(P_{\mathrm E})+o(A^{3/2}).
 \label{eq:individualflatvolume}
\end{equation}
For the separate intrinsic path~\eqref{eq:anglepath}, with fixed
$0<\mu_i<1$, $\sum_i\mu_i=2$, and $d=n-3$, the moduli volume
satisfies
\[
 V_{\WP}(A)=2^{-d}A^d V_{\mathrm{CH}}(\mu)+O(A^{d+1})
\]
by~\eqref{eq:WPCHvolumeasymptotic} in its final chamber. The
normalizations are those of Section~\ref{sec:flatlimit}.
The different powers reflect different dimensions. Flattened
Euclidean limits and changes of polyhedral combinatorics require
additional care; they are not covered by the nondegeneracy
hypothesis in~\eqref{eq:individualflatvolume}.

\subsection{An explicit small-polyhedron expansion}

A controlled model gives more than the leading power. Let $Q$ be a
bounded nondegenerate convex Euclidean polyhedron, and place $P_r=rQ$ in
the curvature-$-1$ Klein ball. Its Euclidean planar faces are
hyperbolic geodesic faces. The Klein metric is
\[
 g_x=\frac{I}{1-|x|^2}
       +\frac{xx^{\mathsf T}}{(1-|x|^2)^2}.
\]
Its two tangential eigenvalues are $(1-|x|^2)^{-1}$ and its
radial eigenvalue is $(1-|x|^2)^{-2}$, so
\[
 \dd\operatorname{Vol}_{\HH^3}=(1-|x|^2)^{-2}\,\dd^3x.
\]
Consequently, whenever $r\sup_Q|y|<1$,
\begin{align}
 W(rQ)
 &=r^3\int_Q(1-r^2|y|^2)^{-2}\,\dd^3y\notag\\
 &=\sum_{k=0}^{\infty}(k+1)r^{3+2k}
                         \int_Q|y|^{2k}\,\dd^3y\notag\\
 &=r^3 W_{\mathrm E}(Q)
      +2r^5\int_Q|y|^2\,\dd^3y
      +3r^7\int_Q|y|^4\,\dd^3y+\cdots.
 \label{eq:Kleinvolumeexpansion}
\end{align}
The series converges uniformly on $Q$ for each
$rR<1$, where $R=\sup_Q|y|$. Its nonnegative remainder after
$k=m$ is bounded by
\[
 0\le W(rQ)-\sum_{k=0}^m(k+1)r^{3+2k}
                 \int_Q|y|^{2k}\,\dd^3y
 \le r^3W_{\mathrm E}(Q)
 \frac{t^{m+1}\{m+2-(m+1)t\}}{(1-t)^2},\qquad t=r^2R^2.
\]
This follows by summing the scalar tail of $(1-t)^{-2}$.
It bounds truncation of this geometric series, not a quantum
asymptotic error. The expansion is exact for the specified
homothetic Klein path. It is not generally the fixed-conformal
angle path~\eqref{eq:anglepath}. Nor does it imply a Taylor expansion
of this form at ideal, flattening, or otherwise uncontrolled
degenerations. Unlike~\eqref{eq:exactWPvolumeflat}, it is an infinite
series: the finite polynomial for the moduli volume comes from a
quadratic cohomology class and a top intersection.

\subsection{Schl\"afli, doubling, and degeneration/regeneration}

For a compact hyperbolic polyhedron with interior dihedral angles
$\alpha_e$ and edge lengths $\ell_e$, Schl\"afli's formula is
\begin{equation}
 \dd W=-\frac12\sum_e\ell_e\,\dd\alpha_e.
 \label{eq:Schlafli}
\end{equation}
For the Klein family above, write
\[
 \ell_e(r)=rL_e+O(r^3),\qquad
 \alpha_e(r)=\alpha_e^{\mathrm E}+r^2b_e+O(r^4).
\]
These parity expansions follow directly from the Klein model.
For example, write two adjacent face planes of $Q$ as
$n_f\cdot y=h_f$ and $n_g\cdot y=h_g$, with unit outward normals.
Their unit Lorentzian normals at scale $r$ are
$(rh_f,n_f)/\sqrt{1-r^2h_f^2}$ and
$(rh_g,n_g)/\sqrt{1-r^2h_g^2}$, for signature $(-,+,+,+)$.
Thus, with $c_e=\cos\alpha_e^{\mathrm E}=-n_f\cdot n_g$,
\[
 \cos\alpha_e(r)
 =\frac{c_e+r^2h_fh_g}
        {\sqrt{(1-r^2h_f^2)(1-r^2h_g^2)}},\qquad
 b_e=-\frac{h_fh_g+\tfrac12c_e(h_f^2+h_g^2)}
                {\sin\alpha_e^{\mathrm E}}.
\]
We take actual faces, so adjacent ones are not coplanar and
$0<\alpha_e^{\mathrm E}<\pi$. The length expansion follows by
integrating the Klein metric along the fixed Euclidean edge,
which is a hyperbolic geodesic. The leading term of
\eqref{eq:Kleinvolumeexpansion} and~\eqref{eq:Schlafli} imply
\begin{equation}
 \boxed{\quad\sum_eL_e b_e=-3W_{\mathrm E}(Q).\quad}
 \label{eq:firstcurvaturevolume}
\end{equation}
Indeed, $\dd W/\dd r=3r^2W_{\mathrm E}+O(r^4)$, while
the right side of~\eqref{eq:Schlafli}, divided by $\dd r$, is
$-r^2\sum_eL_e b_e+O(r^4)$. Thus the first curvature correction
to the Euclidean dihedral angles already determines the Euclidean
enclosed volume. The sign reverses if exterior dihedral angles are
used. As a check, for $Q=[-1,1]^3$ every $L_e=2$ and
$b_e=-1$, so the sum is $-24=-3\cdot8$.
The same identity holds with
$\ell_e=A^{1/2}L_e+O(A^{3/2})$ and
$\alpha_e=\alpha_e^{\mathrm E}+A b_e+O(A^2)$ for a fixed-combinatorics
unit-area limiting family whose expansions are differentiable
in $\sqrt A$. It uses $W=A^{3/2}W_{\mathrm E}+O(A^{5/2})$
with the corresponding differentiated remainder; a bare
$o(A^{3/2})$ estimate cannot be differentiated without control.

Doubling a convex polyhedron across its faces produces a cone
three-manifold with singular graph given by its edges. The cone
angle along an edge is $2\alpha_e$, and its volume is $2W(P)$.
A shrinking polyhedron therefore supplies a direct example of
degeneration to a Euclidean cone three-manifold after rescaling.
This is the geometric setting suggested by Hodgson's thesis,
\emph{Degeneration and regeneration of geometric structures on
3-manifolds}~\cite{HodgsonThesis}. Later precise regeneration results
include Porti--Weiss~\cite[Theorem~1.3]{PortiWeiss}, whose input
is a closed orientable Euclidean cone three-manifold with cone
angles at most $\pi$ which is not almost product. Those
hypotheses are not automatic for a doubled convex polyhedron;
in particular, $2\alpha_e^{\mathrm E}$ can exceed $\pi$. No unrestricted
regeneration theorem is being asserted here. The edge cone angles
$2\alpha_e$ are also distinct from the intrinsic vertex openings
$\theta_i$ of the boundary surface.

Rivin's \emph{Volumes of degenerating polyhedra---on a conjecture
of J.~W. Milnor} is directly relevant~\cite{RivinDegeneration}.
It studies continuous extension of volume toward degenerate angle
data, the Euclidean vanishing boundary, and quantitative control
through Schl\"afli's formula and polar geometry. Its estimates
allow logarithmically growing volume derivatives at degenerations;
continuity does not imply a smooth power series at every boundary
point. The controlled small-polyhedron expansion above gives a
particularly explicit local calculation on the Euclidean side of
this broader degeneration picture.

\subsection{The established tetrahedral quantum connection}

Quantum $6j$-symbols provide a concrete bridge to three-dimensional
geometry. Costantino's hyperbolic-type theorem uses the evaluation
$\operatorname{ev}_n([n]\{6j\})$ after removing a root-of-unity pole;
its geometric value is the volume of the doubled truncated-tetrahedral
block $D(\boldsymbol u)$~\cite[Definitions~3.1--3.2 and Theorem~3.9]{Costantino}.
It is not a formula for an unregularized value at that pole.

For clarity, distinguish two results of Chen--Murakami. Fix a
hyperbolic tetrahedron $T$ with an ideal or ultraideal vertex,
and put $\xi_r=e^{2\pi i/r}$, $r$ odd. Their Theorem~1 gives
\[
 \frac{2\pi}{r}\log|\{6j\}^{\rm RW}_{q=\xi_r^2}|
       \longrightarrow W(T)
\]
for admissible half-integer labels with
$j_{e,r}/r\to(\pi-\theta_e)/(4\pi)$ and with the sum of the
three labels at a chosen ideal or ultraideal vertex at least
$(r-3)/2$. The prefactor requires more: Theorem~2 assumes that
$\eta_e=(\pi-\theta_e)/(2\pi)$ have odd numerators and odd
denominators, takes $r$ through odd common multiples of those
denominators, and uses exactly $j_{e,r}=r\eta_e/2-1/2$. Then
\begin{equation}
 |\{6j\}^{\rm RW}_{q=\xi_r^2}|\sim
 \sqrt2\pi\,r^{-3/2}(-\det G)^{-1/4}
                 \exp\left(\frac{r}{2\pi}W(T)\right),
 \label{eq:quantum6jstructure}
\end{equation}
where $G_{ff}=1$ and $G_{fg}=-\cos\theta_{fg}$ for the incident
face dihedrals~\cite[Theorems~1--2, pp.~1--2]{ChenMurakami}.
The labels are in the source's Racah--Wigner normalization, whose
quantum integer uses $q^{1/2}=\xi_r$; hence $q=\xi_r^2$ here corresponds
to $q_{\rm alg}=\xi_r$ in~\eqref{eq:introqinteger}. Arbitrary
convergent label sequences give the stated leading rate under
Theorem~1, not automatically this refined equivalent.

At the Euclidean end, Roberts proves the Ponzano--Regge asymptotic
for classical $6j$-symbols: Euclidean tetrahedral edge lengths and
dihedral angles enter an oscillatory Regge phase, and Euclidean
volume enters its amplitude~\cite{Roberts}. This is an established
appearance of Euclidean tetrahedral geometry in quantum
asymptotics. It does not automatically identify the classical
large-label limit with a chosen root-of-unity hyperbolic limit.

In particular, the stated Chen--Murakami theorem does not cover
an arbitrary family of compact tetrahedra shrinking to Euclidean
ones. At a Euclidean tetrahedron the angle Gram determinant
vanishes, so the displayed fixed-shape prefactor itself warns of
nonuniformity. A continuous classical volume limit, even with
good quantitative estimates, does not by itself control the
quantum Hessian, the prefactor, or the error term in a joint limit.

\subsection{A proposed joint scale and a model integral}

Suppose a suitable quantum object, in a specified geometric branch
and choice of variables, has an exponential contribution
$\exp(W(P_A)/\hbar)$. Equation~\eqref{eq:individualflatvolume}
then singles out the dimensionless ratio
\begin{equation}
 \lambda=\frac{A^{3/2}}{\hbar},\qquad
 A\asymp\hbar^{2/3}.
 \label{eq:jointscale}
\end{equation}
This is the regime in which geometric degeneration competes
with the semiclassical limit. At this stage it is a scaling prediction;
Section~\ref{sec:quantumAiry} proves an additive uniform Airy
statement for the pole-cleared regular $b$-$6j$ integral and fixed
bounded real transition parameters. In that convention the
volume-action parameter is $\hbar_{\mathrm{vol}}=\pi b^2$,
whereas its fold parameter is $h=2\pi b^2$; see
\eqref{eq:TVvolumehbar}. No claim for fixed nonzero volume or
unbounded transition parameter follows. A different
polarization may involve a Regge or Legendre-transformed action,
so the quantum object and phase must be specified first.

There is a useful classical model making the potential role of
the WP--Thurston limit explicit. Fix the weights and intrinsic
path~\eqref{eq:anglepath}. Let $K\subset\M_{0,n}$ be a compact
measurable subset and assume the corresponding realizations obey
$A^{-3/2}W(P_A(X))\to W_{\mathrm E}(P_{\mathrm E}(X))$ uniformly
on $K$, with a continuous nondegenerate Euclidean limiting family.
Define
\begin{equation}
 \mathcal Z_K(A,\hbar)=\int_K
       \exp\left(\frac{W(P_A(X))}{\hbar}\right)
       \frac{\omega_{\WP,A}^d}{d!}.
 \label{eq:modelpartition}
\end{equation}
This is a constructed geometric integral, not an identification
with a colored-Jones or other known quantum invariant. Take positive $A,\hbar\to0$ with
$A^{3/2}/\hbar\to\lambda_0\in(0,\infty)$. The exponential
then converges uniformly on $K$. By~\eqref{eq:localWPCHlimit},
$2A^{-1}\omega_{\WP,A}\to\omega_{\mathrm{CH},\mu}$ locally
uniformly in form coefficients; taking its $d$th exterior power
preserves this convergence. A finite coordinate cover of $K$
therefore gives a common integrable bound and proves
\begin{equation}
 2^d A^{-d}\mathcal Z_K(A,\hbar)\longrightarrow
 \int_K e^{\lambda_0 W_{\mathrm E}(P_{\mathrm E}(X))}
                 \frac{\omega_{\mathrm{CH},\mu}^d}{d!}.
 \label{eq:modelpartitionlimit}
\end{equation}
Extending this formula to the full noncompact space requires
 domination and control of collision and flattened strata.
 Identifying its measure and phase with those of an actual quantum
 state integral is a further mathematical problem.

This separates the two potentially compatible roles: enclosed
three-dimensional volume supplies an action, while the WP form
supplies a measure on boundary shapes. Their Euclidean limits
are respectively Euclidean polyhedral volume and the
Thurston--Deligne--Mostow measure. Establishing that both occur
together in a particular quantum construction would provide the
substantive connection to the volume conjecture.
